% Титульный лист сборника «Решение задач по теории вероятностей ЕГЭ» (иллюстрация: таблица исходов двух бросков и игральная кость).
\begin{titlepage}
\thispagestyle{empty}
\begin{tikzpicture}[remember picture,overlay]
  % фон
  \shade[top color=titleblue!95!black, bottom color=accent!85!black]
    (current page.south west) rectangle (current page.north east);
  % тонкая сетка
  \begin{scope}[opacity=0.07]
    \foreach \x in {0,0.8,...,21.6} \draw[white,line width=0.3pt]
      ($(current page.south west)+(\x,0)$) -- ++(0,29.7);
    \foreach \y in {0,0.8,...,30} \draw[white,line width=0.3pt]
      ($(current page.south west)+(0,\y)$) -- ++(21,0);
  \end{scope}
  % верхняя плашка
  \node[anchor=north,text=white!80,font=\small\addfontfeatures{LetterSpace=12}]
    at ($(current page.north)+(0,-2.2)$) {ЕДИНЫЙ ГОСУДАРСТВЕННЫЙ ЭКЗАМЕН · МАТЕМАТИКА · ПРОФИЛЬ};
  \draw[white!60,line width=0.6pt] ($(current page.north)+(-4,-2.9)$) -- ++(8,0);

  % заголовок
  \node[anchor=north,text=white,align=center,font=\fontsize{34}{40}\selectfont\bfseries]
    at ($(current page.north)+(0,-4.0)$) {Решение задач\\по теории\\вероятностей ЕГЭ};
  \node[anchor=north,text=white!85,align=center,font=\large]
    at ($(current page.north)+(0,-9.3)$) {Задача № 5 · \TaskCountText\ с подробными решениями};

  % ===== иллюстрация: таблица исходов двух бросков и игральная кость =====
  \begin{scope}[shift={($(current page.center)+(-5.6,-6.2)$)}]
    \fill[white,opacity=0.10,rounded corners=10pt] (-0.8,-0.9) rectangle (12.0,8.4);
    \draw[white,opacity=0.35,rounded corners=10pt,line width=0.8pt] (-0.8,-0.9) rectangle (12.0,8.4);
    % сетка 6 × 6
    \foreach \i in {0,...,6} {
      \draw[white,opacity=0.55,line width=0.8pt] (0,\i) -- (6,\i);
      \draw[white,opacity=0.55,line width=0.8pt] (\i,0) -- (\i,6);
    }
    \foreach \i in {1,...,6} {
      \node[white,font=\small] at (\i-0.5,6.4) {\i};
      \node[white,font=\small] at (-0.4,6.5-\i) {\i};
    }
    % клетки с суммой больше 7
    \foreach \x in {1,...,6} \foreach \y in {1,...,6} {
      \pgfmathtruncatemacro{\s}{\x+\y}
      \ifnum\s>7 \node[yellow!80!orange,font=\large\bfseries] at (\y-0.5,6.5-\x) {+}; \fi
    }
    % игральная кость
    \begin{scope}[shift={(7.6,3.2)}]
      \fill[white,rounded corners=8pt] (0,0) rectangle (3,3);
      \foreach \p in {(0.7,0.7),(0.7,1.5),(0.7,2.3),(2.3,0.7),(2.3,1.5),(2.3,2.3)}
        \fill[titleblue] \p circle (0.24);
    \end{scope}
    \node[text=white,font=\Large,anchor=west] at (7.3,1.4) {$P = \dfrac{m}{n}$};
    \node[text=white,font=\large,anchor=west] at (7.3,7.35) {$\dfrac{15}{36}\approx 0{,}42$};
  \end{scope}

  % подвал
  \node[anchor=south,text=white!85,align=center,font=\normalsize]
    at ($(current page.south)+(0,3.0)$) {Кости и монеты · полная вероятность · независимые события · сочетания};
  \draw[white!60,line width=0.6pt] ($(current page.south)+(-4,2.6)$) -- ++(8,0);
  \node[anchor=south,text=white!70,font=\small\addfontfeatures{LetterSpace=10}]
    at ($(current page.south)+(0,1.7)$) {\CollectionYear};
\end{tikzpicture}
\end{titlepage}
